
%*****************************************
\chapter{Conclusion}
\label{ch:closure}
%*****************************************
This work addressed the challenge of improving robustness and calibration in neural networks for EEG-based seizure classification, with a particular focus on real-world distribution shifts and noise corruptions. A new benchmark, TUSZ-C, was introduced, systematically incorporating clinically inspired corruptions such as Gaussian noise, signal dropout, spectral augmentation, and channel dropout. This benchmark enabled controlled robustness evaluation under settings that more accurately reflect the conditions faced in clinical neurophysiology. 

A comprehensive empirical study was conducted, evaluating multiple strategies to enhance model reliability: Multi-Input Multi-Output (MIMO) architectures, Manifold Mixup, and their combination. Performance was measured not only in terms of classification accuracy (F1-score) but also calibration quality (ECE), emphasizing the importance of confidence alignment in high-stakes decision-making. Key findings revealed that MIMO consistently offered superior robustness and calibration across corruptions, while Manifold Mixup provided improvements primarily under additive and frequency-based distortions. However, combining both methods did not yield additive benefits, underscoring the complexity of interacting regularization mechanisms. In addition, test-time adaptation via entropy minimization proved to be a powerful post-hoc calibration method, especially for models suffering from overconfidence under distribution shift.

Despite its contributions, this work has several limitations. First, the evaluation was limited to a single dataset (TUSZ), and the generalizability of the findings to other EEG corpora or clinical tasks (e.g., multi-class classification, long-term monitoring) remains to be tested. Second, while five corruption types were introduced, further corruptions, including biologically plausible artifacts such as eye blinks or muscle noise, could offer deeper insight into real-world robustness. Third, the interplay between model architecture, augmentation strategy, and data imbalance was not fully disentangled; future work could involve systematic ablation studies or calibration-aware training objectives. Finally, while TTA was effective, it was only tested using entropy minimization. Exploring alternative unsupervised adaptation strategies such as test-time feature alignment, self-supervised pretext tasks, or memory-based adaptation, could further enhance post-hoc robustness without incurring high deployment costs.